\documentclass{article}

\def\CourseCode{ECEN-662}
\def\CourseName{Estimation \& Detection Theory}
\def\Instructor{Dr. Tie Liu}
<<<<<<< HEAD
\def\Semester{Spring 2026}
\def\Student{Brody Jordan}
\def\University{Texas A\&M University}
\def\Cover{figures/_cover.pdf}
=======
\def\Semester{Spring 2026}
>>>>>>> 249e4c317b2a4234368ca0c3c75e714883e8075a

\usepackage[
    AssignmentName={Homework 6},
    DueDate={March 25, 2026 at 11:59 PM}
]{assignment}

\usepackage{cancel}

\begin{document}
\MakeAssignmentTitle

% ------ PROBLEM 1 ------ %
\begin{homeworkProblem}

  Consider a Gaussian Mixture Model (GMM) with two well-separated modes:
  \[
    p(\theta) = 0.5 \cdot \Nc(\theta;-3, 1) + 0.5 \cdot \Nc(\theta;3, 1),
  \]
  where $\Nc(\theta;\mu,\sigma^2)$ is a Gaussian density with mean $\mu$ and variance $\sigma^2$. Our goal is to approximate the above GMM using the variational
  distribution:
  \[
    q_\lambda(\theta) = \Nc(\theta; \mu, \sigma^2),
  \]
  where the variational parameter $\lambda = (\mu, \sigma)$.
  \begin{enumerate}[label=\roman*), topsep=5pt, itemsep=5pt]
    \item Derive the explicit update rule for both $\mu$ and $\sigma$ using the re-parameterization trick.
    \item Write a computer program to find the optimal variational distribution for the
          above GMM model. Try multiple initial values for $\mu$ and $\sigma$ and describe how the value of
          $\mu$ and $\sigma$ change over the iterations.
    \item Repeat Part ii) for the following updated GMM model:
          \[
            p(\theta) = 0.5 \cdot \Nc(\theta;-3, 4) + 0.5 \cdot \Nc(\theta;3, 0.25).
          \]
  \end{enumerate}

  \solution

  \textbf{Part (i):} ELBO:
  \[
    \Lc(\mu, \sigma) = \E_{q_{\lambda}} [\log q_\lambda(\theta) - \log p (\theta)]
  \]
  So
  \[
    \Lc(\mu, \sigma) = \E_{\epsilon \sim \Nc(0, 1)} [\log q_\lambda(\mu + \sigma \epsilon) - \log p (\mu + \sigma \epsilon)]
  \]
  Which becomes
  \[
    \Lc(\mu, \sigma) = - \frac{1}{2} \log (2\pi \sigma^2) - \E_{\epsilon \sim \Nc(0, 1)} [\log p (\mu + \sigma \epsilon)]
  \]
  Now taking the gradients:
  \[
    \nabla_\mu \Lc(\mu, \sigma) = - \E_{\epsilon \sim \Nc(0, 1)} [\nabla_\theta \log p (\theta)]
  \]
  \[
    \nabla_\sigma \Lc(\mu, \sigma) = - \frac{1}{\sigma} - \E_{\epsilon \sim \Nc(0, 1)} [\epsilon \cdot \nabla_\theta \log p (\theta)]
  \]
  Thus we update $\mu$ an $\sigma$ as follows:
  \begin{align*}
    \mu    & \leftarrow \mu + \alpha \nabla_\theta \log p (\mu + \sigma \epsilon)                                                     \\
    \sigma & \leftarrow \sigma + \alpha \left( \frac{1}{\sigma} + \epsilon \cdot \nabla_\theta \log p (\mu + \sigma \epsilon) \right)
  \end{align*}

  \textbf{Part (ii):}

  \begin{table}[hp!]
    \centering
    \begin{tabular}{c|c|c}
      Initial $\mu$ & Initial $\sigma$ & Final $(\mu, \sigma)$ \\
      \hline
      -3            & 1                & (-3.04, 1.01)         \\
      -3            & 3                & (-3.02, 1.01)         \\
      0             & 1                & (-0.13, 2.74)         \\
      0             & 3                & (0.16, 2.84)          \\
      3             & 1                & (3.06, 1.00)          \\
      3             & 3                & (2.96, 1.03)          \\
    \end{tabular}
  \end{table}
  \pagebreak

  \textbf{Part (iii):} Rerunning but now with
  \[
    p(\theta) = 0.5 \cdot \Nc(\theta;-3, 4) + 0.5 \cdot \Nc(\theta;3, 0.25)
  \]
  yields

  \begin{table}[hp!]
    \centering
    \begin{tabular}{c|c|c}
      Initial $\mu$ & Initial $\sigma$ & Final $(\mu, \sigma)$ \\
      \hline
      -3            & 1                & (-2.04, 2.79)         \\
      -3            & 3                & (-1.96, 2.62)         \\
      0             & 1                & (-2.06, 2.82)         \\
      0             & 3                & (-2.09, 2.62)         \\
      3             & 1                & (3.01, 0.51)          \\
      3             & 3                & (3.01, 0.51)          \\
    \end{tabular}
  \end{table}



\end{homeworkProblem}

\end{document}

